\clearpage
\phantomsection% 
\addcontentsline{toc}{chapter}{ABSTRAK}
\thispagestyle{fancy}

\begin{center}
	\large \bfseries \MakeUppercase{Abstrak}\\
%	\normalsize \normalfont {\thetitle}\\
%	\normalsize \normalfont {\theauthor}\\
	\bigskip
	
	\normalsize \normalfont \justifying \singlespacing
	Ketidakseimbangan distribusi kelas merupakan permasalahan yang sering dihadapi dalam klasifikasi \textit{machine learning}, menyebabkan model cenderung bias terhadap kelas mayoritas dan gagal mendeteksi kelas minoritas. HOUM (\textit{Hybrid Oversampling and Undersampling Method}) merupakan metode hibrida yang menggabungkan Safe-Level SMOTE untuk \textit{oversampling} dan SVM untuk \textit{undersampling} secara iteratif. Namun, penelitian sebelumnya belum mengeksplorasi pengaruh variasi parameter regularisasi $C$ dan jenis kernel SVM, serta belum membandingkan Safe-Level SMOTE dengan teknik \textit{oversampling} lain seperti ADASYN. Penelitian ini bertujuan untuk menganalisis pengaruh variasi teknik \textit{oversampling} (Safe-Level SMOTE vs ADASYN) dan konfigurasi SVM (parameter $C$: 0,5; 1; 1,5; 2; 2,5; 3 serta kernel Linear, RBF, Polynomial) pada metode HOUM terhadap kinerja klasifikasi data tidak seimbang. Empat dataset \textit{benchmark} digunakan, yaitu Diabetes, Haberman, Blood Transfusion, dan Glass5 dengan variasi rasio ketidakseimbangan dari 1:1,87 hingga 1:22,78. Data dibagi 80:20 secara \textit{stratified}, diseimbangkan dengan 18 kombinasi konfigurasi HOUM, dan diklasifikasi menggunakan \textit{Random Forest}. Hasil menunjukkan HOUM-SLSMOTE konsisten menghasilkan \textit{recall} lebih tinggi pada Diabetes (0,6481), Haberman (0,3750), dan Transfusion (0,7500), sedangkan HOUM-ADASYN unggul pada \textit{precision} di Diabetes (0,7234) dan Glass5 (0,6667). Pada Glass5 dengan ketidakseimbangan ekstrem, HOUM-ADASYN mencapai \textit{recall} 1,0000, ROC-AUC 1,0000, dan G-Mean 0,9877. Nilai $C=0,5$ optimal pada Diabetes, Haberman, dan Glass5, sedangkan $C=2$ lebih sesuai pada Transfusion. Kernel RBF unggul pada dataset berdimensi tinggi (Diabetes), kernel Linear pada dimensi rendah (Haberman, Transfusion), dan pada Glass5 seluruh kernel menghasilkan metrik identik pada HOUM-ADASYN. Konfigurasi terbaik: HOUM-ADASYN RBF $C=0,5$ (Diabetes, G-Mean 0,7401), HOUM-SLSMOTE Linear $C=0,5$ (Haberman, G-Mean 0,4778), HOUM-SLSMOTE Linear $C=2$ (Transfusion, G-Mean 0,6977), dan HOUM-ADASYN Linear $C=0,5$ (Glass5, G-Mean 0,9877). Kesimpulannya, pemilihan teknik \textit{oversampling}, kernel SVM, dan nilai $C$ harus disesuaikan dengan karakteristik dataset karena terdapat \textit{trade-off} antara \textit{recall} dan \textit{precision}. HOUM terbukti efektif meningkatkan deteksi kelas minoritas, termasuk pada kasus ketidakseimbangan ekstrem seperti Glass5.

    \vspace{20pt}
	\raggedright \textbf{Kata Kunci: HOUM, Oversampling, Undersampling, ADASYN, Safe-Level SMOTE, SVM, Imbalanced Data}
	
	\vfill
	
\end{center}
\clearpage